\section{Convex duality and sharp operator bounds}
\label{sec:geometry-operators}

The Mahler and Crouzeix families are useful tests of the release's proposed
mathematical mechanisms: both seek sharp universal inequalities, but the
obstruction is not simply estimating another integral more accurately.
For Mahler it is retaining enough information about polar geometry to
identify the extremizers; for complete Crouzeix it is preserving
noncommutative multiplication under arbitrary matrix amplification.
The discussion below follows the released proof passages and identifies
the transitions on which the claimed advances depend.

\subsection{Mahler: boundary feasibility and Gaussian spectral comparison}
\label{subsec:go-mahler}

A volume product measures how a convex body and its polar constrain one
another; it is unchanged by invertible linear changes of coordinates.
Sharp lower bounds therefore identify geometric shapes rather than a
preferred coordinate system. For a compact convex body
\(K\subset\mathbb R^n\) with nonempty interior,
write \(K^\circ=\{y:\langle y,x\rangle\le1\ (x\in K)\}\) when the origin is
interior. The symmetric manuscript claims, for every \(n\ge1\),
\[
 K=-K
 \quad\Longrightarrow\quad
 |K|\,|K^\circ|\ge\frac{4^n}{n!},
\]
with equality precisely for invertible linear images of Hanner bodies:
bodies generated from centered intervals by Cartesian products and
convex-hull joins. The general manuscript separately claims
\[
 \inf_{z\in\operatorname{int}K}|K|\,|(K-z)^\circ|
 \ge\frac{(n+1)^{n+1}}{(n!)^2},
 \qquad
 \text{equality if and only if \(K\) is a simplex}.
\]
Neither statement assumes smoothness or strict convexity
\cite{GOOpenAISymmetricMahler2026,GOOpenAIGeneralMahler2026}.
The constants and extremal classes differ; the general inequality does
not recover the sharp symmetric one by restriction.

The symmetric proof starts with a strip polytope
\(A=\{X:|b_iX|\le1\}\), using nonzero, pairwise nonproportional rows
that span the dual space after removing redundant parallel strips.
Its unusual connecting device is a conformal lens. A disk map \(F\)
has boundary law
\[
 F(e^{i\Theta})\overset{\rm law}{=}\varepsilon\lambda(T)+iT,
 \qquad T\sim\operatorname{Unif}[-1,1],
\]
with an independent fair sign \(\varepsilon\), and a strictly concave
half-width \(\lambda\). The manuscript explicitly obtains this map by
rotating Gross's uniform-distribution example
\cite[Section~3.2]{GOGross2019}. Thus the conformal lens is a prior
construction; the proposed advance is its use to convert complex
feasibility probabilities into sharp real volume information.

For each independent row basis \(I\), a boundary sample determines a
point \(Z_I=B_I^{-1}(F(e^{i\theta_i}))_{i\in I}\). Let \(P_I\) be the
probability that every remaining strip coordinate belongs to the closed
lens. The boundary-feasibility lemma claims \(\sum_I P_I\ge1\).
Its proof applies a holomorphic isolated-zero mass estimate to the
high powers \(g(b_i z)^k\), with \(g=F^{-1}\). The lower mass grows as
\(k^n\); the substitution \(u_i=\rho_i^{1/k}e^{i\theta_i}\) removes
that factor and upper Fatou passes to closed boundary feasibility.
This mass estimate belongs to the generalized Lelong-number framework
\cite[Sections~3 and~5]{GODemailly1993}; complex-analytic approaches
to reverse volume-product bounds also precede the release
\cite{GOBerndtsson2021}. The noteworthy proposed step is the
dimension-independent feasibility inequality with its exact normalization.

The next transition is more geometric. For fixed \(X\), feasible signed
bases give simplices in \(A^\circ\). Additional boundary ties are shown
to have probability zero, and a supporting-functional argument proves
that the simplex interiors are disjoint. The resulting identity is
\[
 \sum_I P_I
 =\frac{n!}{4^n}\int_A|\Sigma_X|\,dX
 \le\frac{n!}{4^n}|A|\,|A^\circ|,
\]
where \(\Sigma_X\) is their union. This exposes where the sharp constant
enters. Approximation of arbitrary symmetric bodies by polars of inner
polytopes then extends the bound. Compared with the established
three-dimensional theorem and equality classification of Iriyeh--Shibata
\cite{GOIriyehShibata2020}, the claimed scope is all dimensions through
an analytic feasibility identity rather than a dimension-specific
partition argument.
All-dimensional nonsharp lower bounds also precede the release, including
Kuperberg's indefinite-geometric and linking-integral approach
\cite{GOKuperberg2008}. The proposed qualitative gain is the conjectured
sharp endpoint together with full rigidity, rather than a first
all-dimensional reverse volume-product estimate.

\begin{example}[A geometric transfer calibrated on the square]
Take $A=[-1,1]^2$ with rows $e_1,e_2$. There is one indexed basis,
its feasibility probability is $1$, and the four signed polar triangles
fill the diamond $A^\circ$, of area $2$. Hence
$(2!/4^2)\int_A|\Sigma_X|\,dX=(1/8)\cdot4\cdot2=1$.
Representing the same square by rows $e_1,e_1,e_2$ instead gives two
indexed bases, $\{1,3\}$ and $\{2,3\}$, each with probability $1$.
Their geometric union remains the same diamond. Thus the probability
sum is $2$ while the normalized union volume is $1$.
Removing proportional rows is consequently part of the sharp geometric
transfer: preserving the body alone does not preserve the indexed
probability accounting. This checks the distinction between the
feasibility inequality and the disjoint-simplex identity under their
respective hypotheses.
\end{example}

Equality requires a separate mechanism. After adjoining an
\(\ell_1\)-segment, vanishing integrated uncovered volume is tested on
product neighborhoods of a prescribed interior point and direction.
Compactness gives a single feasibility witness, even when limiting
row bases become dependent. The lens endpoint asymptotic
\(\lambda(1-\eta)\sim(2/\pi)\eta\log(1/\eta)\) converts minimum
representation costs into the existence of metric medians for every
triple in the polar norm. The final three-ball classification is the
classical Hansen--Lima theorem
\cite[Corollary~7.4]{GOHansenLima1981}. Hence the proposed rigidity
advance lies in deriving that classical norm property from sharp
volume equality. Checking the mass normalization, disjointness identity,
and integral-to-universal-witness passage is more decisive than checking
the concluding classification alone.

The general proof uses a substantially different structure. Lift the
body to a closed, convex, solid, pointed cone \(C\subset\mathbb R^m\),
\(m=n+1\), and use its positive dual \(D\). With interior vectors
\(U\in C,V\in D\) normalized by \(\langle U,V\rangle=m\), define
\(\chi_C(V)=\int_Ce^{-\langle V,x\rangle}\,dx\).
The desired statement becomes
\(\chi_C(V)\chi_D(U)\ge1\). This cone--Laplace correspondence is
already explicit in Klartag's work
\cite[Lemma~2.1 and Corollary~4.1]{GOKlartag2018}.
The new claim concerns the estimate after that reduction, rather than
the reduction itself.

Biased Gaussian projections onto the two cones supply cone-valued
maps and projection derivatives \(0\preceq P_z\preceq I\).
A simultaneous normalization chooses both linear coordinates and
covariance feedback through a fixed-point argument, including a
noncommuting matrix quadratic. Gaussian smoothing, injectivity, change
of variables, and log-determinant concavity lower-bound the normalized
logarithm of the Laplace product by \(-\mathcal E\). Hermite-degree
separation and inverse Ornstein--Uhlenbeck covariance representations
then compare the entire layer field with spectral thresholds of a
linear Gaussian matrix. The covariance machinery has primary precedents
in Houdr\'e--Privault
\cite[Proposition~2.1]{GOHoudrePrivault2002}.

The purpose of \(\mathcal E\) is to collect the remaining terms in
this lower bound. Proving \(\mathcal E\le0\) makes the logarithmic
lower bound nonnegative and hence gives the desired product at least
one. The hard step is therefore controlling these terms together,
including those caused by noncommuting matrices.

A particularly distinctive ingredient is the tuned auxiliary curve:
one component rotates with phase \(4.6\operatorname{arsinh}(x/3.6)\);
another is chosen so the scalar profile plus the curve's squared norm
is affine in that component. A strict inequality for every nontrivial real segment
absorbs divided-difference, layer, and commutator errors and establishes
\(\mathcal E\le0\). Equality forces commuting deterministic matrices,
then diagonal cone-projection derivatives and an orthant cone,
yielding a simplex section. This is global rigidity, beyond the prior
local simplex minimality and stability theorem
\cite{GOKimReisner2011}. The delicate proposed advance is the
simultaneous matrix normalization and all-segment absorption inequality,
not Gaussian integration by parts in isolation.
The June 2026 Chen--Li--Xi--Xu preprint already claims the general
three-dimensional inequality and equality classification using shadow
flows \cite{GOChenLiXiXu2026}; the release's proposed extension is
dimension-independent and follows a different proof architecture.

The scalar appendices provide exact rational node enclosures together
with analytic interpolation and tail bounds. Their finite
arithmetic certificates and the whole-real-line interpolation and tail
arguments supply complementary checks: the former enclose specified
nodes, while the latter connect those enclosures to the all-segment
estimate used in the noncommuting matrix comparison.

\subsection{Crouzeix: retaining order under matrix amplification}
\label{subsec:go-crouzeix}

For \(A\in M_n(\mathbb C)\), let
\(W(A)=\{x^*Ax:\|x\|_2=1\}\).
The direct manuscript claims, for all positive \(n,m\), all degrees
\(d\ge0\), and arbitrary \(B_k\in M_m(\mathbb C)\),
\[
 \left\|\sum_{k=0}^d A^k\otimes B_k\right\|_2
 \le2\max_{z\in W(A)}
       \left\|\sum_{k=0}^d z^kB_k\right\|_2 .
\]
The coefficient dimension is independent of the base dimension.
Scalar polynomials correspond to \(m=1\); allowing every \(m\) with
the same constant is what makes the assertion \emph{complete}. It
controls matrix-valued functions of \(A\), whose coefficient products
need not commute. Points and segments are included \cite{GOOpenAIDirectCrouzeix2026}.
For a concrete sharpness control, take
\[
 A=\begin{pmatrix}0&2\\0&0\end{pmatrix},\qquad p(z)=z.
\]
For a unit vector \((x_1,x_2)\),
\(x^*Ax=2\overline{x_1}x_2\); its possible values fill the closed
unit disk, while \(\|A\|_2=2\). Thus constant two is already necessary
for scalar coefficients. Sharpness and matrix amplification are
separate parts of the proposed theorem. The established complete bound
\(1+\sqrt2\) supplies the unconditional comparison
\cite{GOCrouzeixPalencia2017}. Positive double-layer representations
and convex spectral-set methods are older ingredients
\cite{GODelyonDelyon1999,GOBadeaCrouzeixDelyon2006}.

The direct proof works first on an analytic convex enclosing domain
\(\Omega\), with the matrix-valued analytic test normalized by
\(\sup_{z\in\Omega}\|F(z)\|\le1\).
An exterior conformal coordinate produces Faber coefficients and a
nonnegative boundary kernel having unit mass in both angular variables.
The associated Fourier map is an \(L^2\) contraction: the negative
Fourier part is controlled by the positive part. This yields a
Hilbert--Schmidt coefficient estimate with weight one for the constant
coefficient and weight two for all positive degrees.

The decisive algebraic passage selects a top singular pair \(x,y\)
of \(F[A]\), compresses the positive Hermitian resolvent field, and tests
the coefficient functional separately on \(FG\) and \(GF\).
These ordered tests bound both diagonal blocks using reciprocal Fourier
weights. The off-diagonal block retains the constant term and has a
matching lower bound. Positivity of the resulting \(2\)-by-\(2\) block
forces \(\|F[A]\|\le2\). The final block inequality is elementary;
the substantive proposed improvement is obtaining its compatible
diagonal and cross estimates uniformly in coefficient size.

This scope matters for current positioning. Lorist--Schwenninger's
August 2026 manuscript claims scalar constant two using uniformly
bounded perturbations of \(2\)-dilations and all powers \(f^j\).
Its Section~3(iv) explicitly identifies the commutativity obstruction
to direct matrix amplification
\cite{GOLoristSchwenninger2026}. The September version of
\AA hag--Czy\.z--Virtanen claims the complete inequality for base
matrices of order at most three
\cite{GOAhagCzyzVirtanen2026}. Those preceding claims preclude
describing every constant-two statement as new here; they do not by
themselves supply the release's arbitrary-dimensional complete theorem.

Two earlier scalar preprints are also close at the level of mechanism.
Jin's August version uses positive-real completion and ordered weighted
Gramians; Luo's August version uses top singular vectors, positive double
layers, and Fourier or defect-Gram boundary certificates
\cite{GOJin2026,GOLuo2026}. Both claim scalar endpoints and expressly
distinguish their scope from complete matrix-coefficient estimates.
They narrow the novelty judgement: positive kernels, Gram certificates,
and boundary cancellation already occur in preceding proposed routes.
The release's distinctive target is the compatibility of its ordered
estimates under arbitrary matrix amplification.

The structural companion formulates the same sharp target through an
attained minimum similarity: for a finite matrix whose numerical range
lies inside a bounded convex domain with regular analytic Jordan
boundary, the disk-coordinate operator is similar to a contraction
with condition number at most two. A common continuous positive
boundary density of mass \(I\) represents all independently chosen
matrix-valued analytic tests \cite{GOOpenAIStructuralCrouzeix2026}.
Paulsen's completely bounded similarity theorem already supplies the
matching condition number and attained similarity for a unital
completely bounded homomorphism \cite{GOPaulsen1984}.
The companion also derives its common density before establishing
the sharp condition-number bound, using the classical contraction
resolvent representation and crediting Delyon--Delyon.
Similarity, attainment, and density existence are therefore inherited
parts of this architecture. The candidate advance is the bound on the
optimal condition number through the ordered density-comparison
lemma; separate left and right
multiplication estimates again preserve amplification. The companion
also claims extension to bounded operators on arbitrary complex
Hilbert spaces and functions holomorphic near the closed numerical
range, extending the target beyond finite-dimensional matrices.

Across both families, the proposed sharp constants depend on a
compatibility that the classical ingredients do not separately supply:
complex mass must match the disjoint polar simplices, scalar enclosures
must control a global matrix comparison, and ordered Fourier estimates
must survive amplification. This locates the mathematical contribution
more precisely than attributing novelty to the use of positivity itself.
The supplied formal targets describe the corresponding sharp
statements.\footnote{The scope documents link comparator templates
\cite{GOOpenAIScope0872026,GOOpenAIScope3252026}. In family 087, a
sentence excluding nonsymmetric Mahler conflicts with the linked
general-Mahler target; the target declares the general inequality and
simplex equality. The comparison follows the target statement.}
